% Chapter 15 — Testability and the Executable-Check Regime (WILSON). Budget 5.5 pp.
% Discharges Constitution v1.1 §13 (two layers of correctness) and the §2 Testability commitment.
\section{Testability and the Executable-Check Regime}\label{ch:testability}

This chapter discharges \S 13 of the constitution --- the two layers of correctness ---
and the Testability commitment of \S 2. The claim it makes good is the constitution's: the
ledger is correct because it is built to be tested. Every invariant catalogued in
Chapter~\ref{ch:invariants} is an executable check run over generated products, events, and
histories, never a property defended only in prose; and correctness composes, because the
steps are pure and their composition is fixed --- the map, then the fold --- so verifying
each part verifies the whole. This chapter does not restate the invariants. It makes them
the surface a test regime exercises, states the regime, and turns the specification's three
binding conditions into executable gates rather than intentions.

\subsection{Two layers of correctness}

Correctness rests on two layers, and each is checked by its own means.

\emph{Structural} correctness holds by construction, at the door. An illegal state is not
detected and rejected; it is unrepresentable --- no program writes it. A balance changes
only through a move (Chapter~\ref{ch:objects}); a marker coordinate is unwritable without
its governing agreement reference (Chapter~\ref{ch:collateral}); a knocked note has no
un-knock constructor (Chapter~\ref{ch:objects}, Invariant~\ref{inv:graph-consistency});
conservation is automatic because every move is paired-leg and the system is closed. These
properties hold whatever the contract that proposed the transaction did --- correct,
careless, or malicious --- because the Transaction Executor admits on structure alone and
is the sole writer (Chapter~\ref{ch:machines}). The test obligation for this layer is
negative: no generated event, however adversarial, drives the ledger into a state the type
discipline forbids or the door should refuse.

\emph{Economic} correctness is not guaranteed at admission; it is made checkable by
recomputation. A contract may emit a perfectly well-formed transaction that books the wrong
economics. The check is to re-run the map on the recorded event and compare the result with
what was recorded.

\begin{lstlisting}
-- the economic oracle: re-run the map on the recorded event and compare
prop_recompute :: Recorded (Event, Transaction) -> Property
prop_recompute (e, txRecorded) = contract e (ledgerBefore e) === txRecorded
-- a mismatch is a defect of the CONTRACT, not the ledger; it is repaired by a
-- compensating transaction through the same door -- never by an edit to history.
\end{lstlisting}

Purity is what makes this oracle exist at all: a contract reads only the event and the
record (Chapter~\ref{ch:contracts}), so its recomputation is exact and unconditional. A
mismatch is a defect of the contract, never of the ledger, and it is repaired the way every
change is made --- by a compensating transaction through the same door. The regime-bit
repair of Chapter~\ref{ch:collateral} (a new terms version \emph{plus} the compensating
rebooking of the affected interval) is the worked instance; the valuations of
Chapter~\ref{ch:valuation} are recomputable under the same discipline. The two layers
divide the labour exactly and neither substitutes for the other: structure is proved once
and tested negatively, economics is re-derived per recorded event and tested by comparison,
and a structurally valid transaction can still be economically wrong.

\subsection{The generator universe and the adversarial scheduler}

The enumerations this specification builds on are closed: the lifecycle states and edges of
a product graph, fixed at registration (Chapter~\ref{ch:objects}); the collateral regimes
and marker planes (Chapter~\ref{ch:collateral}); the component dimensions of the datum-kind
registry and the boundary failure regimes W1--W4 (Chapter~\ref{ch:marketdata}); the trigger
classes of a contract (Chapter~\ref{ch:contracts}). Because they are closed and declared,
one generator universe ranges over the whole domain, and its coverage of each enumeration is
itself a checkable fact --- every inhabitant is reachable by the generators, so a property
that never exercises a case is a visible gap, not a silent one.

\begin{lstlisting}
genProduct :: Gen ProductGraph          -- a well-typed graph over the closed node/edge set
genEvent   :: ProductGraph -> Gen Event -- on-menu guards, and off-menu adversarial events
genHistory :: Gen [Transaction]         -- interleaved admitted walks of many products
-- every generator is paired with a shrinker: a violation shrinks to a minimal history,
-- and every counterexample is replayable from its seed (a defect is a reproduction).
\end{lstlisting}

Generation is adversarial by design, because the test environment must be harder than
production. The Event Monitor and the Events Executor are outside the trust boundary and may
duplicate, delay, and reorder firings; integrity owes them nothing
(Chapter~\ref{ch:machines}). The scheduler therefore explores every legal interleaving the
Transaction Executor's total order permits --- duplicated proposals, late firings, events
delivered out of arrival order --- and injects pathological-but-legal behaviour rather than
only the happy path: this is precisely what exercises idempotence and the temporal model
below. One discipline guards the generators against vacuity: a property whose precondition
is never generated passes without witnessing anything. A generated history that never knocks
a pledged note, never re-posts received collateral, never delivers a duplicate is not
evidence; the generators are measured by the states they reach, not the cases they cover.

\subsection{The test surface: the invariant catalogue and the product graph}

The invariants of Chapter~\ref{ch:invariants} are the test surface, in a ramp from the
universal to the domain-specific. At the base sit the properties no history may violate: no
transaction half-applies (atomicity), no balance is created or destroyed off the paired-leg
move (conservation per (unit, coordinate)), the log is append-only and hash-chained, and
every arrow is deterministic. Above them sit the structural invariants --- consistency of
reference, writer discipline --- and above those the safety properties of the collateral
planes: coverage is an invariant, checked at the door on every admission, while sufficiency
is an obligation, lawfully false intraday and bounded by a deadline, and the test regime
must assert the first as an invariant and the second only as a liveness obligation. Stating
sufficiency as an invariant would be a false theorem the generators refute at once, and
insisting on it would forbid lawful states: the two guards are never collapsed, in the tests
least of all.

Carried onto this surface are the four properties of the product-graph consistency invariant
(Invariant~\ref{inv:graph-consistency}, defined in Chapter~\ref{ch:objects}), stated here as
executable checks over generated products and histories:

\begin{lstlisting}
-- for every unit, at every commit, the three interpreters of its graph agree:
prop_watchesAreOutEdges   :: History -> Property  -- (i) watch list == out-edges(current node)
prop_firedMatchesOneEdge  :: History -> Property  -- (ii) fired event -> exactly one edge, lands on its target
prop_actionEqualsRecorded :: History -> Property  -- (iii) emitted tx == edge action applied to recorded inputs
prop_traversalExpiresSibs :: History -> Property  -- (iv) traversal expires exactly the declared siblings
\end{lstlisting}

Property~(i) is checked at every log position: the Event Monitor's live watch set is folded
from the record and compared with the out-edge set of the unit's current node. Property~(ii)
rejects any firing that matches no edge or more than one, or that lands the unit off the
edge's declared target. Property~(iii) is the graph-level face of the economic oracle above:
the transaction the contract emitted must equal the edge's declared action applied to the
recorded inputs. Property~(iv) walks the sibling set and confirms that a traversal expires
those edges and no others. The graph declares watches, state schema, and actions from one
object; these four properties are the statement, made executable, that its three interpreters
never disagree.

\subsection{The three binding conditions as executable gates}

Three conditions the specification relies on are stated here as gates a candidate
implementation must pass, not as intentions it may assert.

\paragraph{(i) Moves are the sole balance mutator.}
The constitution makes the move the only operation that changes a balance
(Chapter~\ref{ch:objects}); the gate makes that mechanically decidable. Applying a
transaction, then applying only its move list, must leave identical balances --- so a
unit-state, position-state, or terms change that touched a balance is caught as a divergence.

\begin{lstlisting}
prop_movesOnly :: Transaction -> Property
prop_movesOnly tx = balances (apply tx l0) === balances (apply (onlyMoves tx) l0)
  where l0 = ledgerBefore tx
-- corollary, tested jointly:  apply (dropMoves tx) l0  leaves every balance fixed.
\end{lstlisting}

\paragraph{(ii) Determinism and replay.}
The clock is the single input not on the record, and the Event Monitor is the single place it
is consulted (Chapter~\ref{ch:machines}); every arrow downstream is a function of the record
alone. Two gates follow. \emph{Idempotence:} a duplicated or late proposal, carrying the same
cause-derived identifier, commits once and produces the identical transaction. \emph{Replay
determinism:} re-running the map on the recorded observations reproduces every recorded
statistic bit-for-bit, because barrier and fixing observation are deterministic --- each
datum crosses the boundary once, as a recorded observation, and is never re-read live
(Chapter~\ref{ch:marketdata}).

\medskip\noindent\textbf{Episode --- the future (T1): a duplicate is inert; a late firing is
identical.}
FUT-IDX settles at recorded IDX closes $1{,}000$, $990$, $1{,}005$ (multiplier $10$);
cumulative variation margin is $+50-100+150 = 100$ (Chapters~\ref{ch:contracts}
and~\ref{ch:valuation}). \emph{The gate.} The day-2 settlement event is delivered twice, and
delivered late, in a scheduler that reorders it against unrelated firings. The contract is a
function of the record, so the late firing emits the identical \USD{100} move; the duplicate
carries the same cause-derived identifier and commits once at the door. \emph{Design point.}
\emph{Timeliness depends on the untrusted machines; integrity does not --- a duplicate is
inert and a late firing identical, and both are properties the scheduler tries and fails to
break.}

\medskip\noindent\textbf{Episode --- the variance swap (T5): regenerate firing~127.}
VS-1 accrues a sum-of-squared-returns statistic in UnitStatus, an integer at scale $10^{-8}$,
written by each firing for the next (Chapter~\ref{ch:homes}). At the mid-life checkpoint the
accrued statistic is $A_{126} = 1{,}805{,}000$, and the previous recorded close is fixing
$126 = 1{,}005$. Firing~127 reads $A_{126}$ from the record, forms the return against the
recorded closes, and writes $A_{127}$.

\begin{lstlisting}
accrue :: Close -> Close -> Accrued -> Accrued        -- named, versioned, pure
prop_replay127 :: Property
prop_replay127 = accrue close127 1005 1805000 === recorded A_127
-- the recorded closes are deterministic; A_126 is a recorded integer; so any party
-- re-running the same map on the same recorded fixings obtains the identical A_127.
\end{lstlisting}

Firing~127 is not trusted to have been right because it ran; it is verified by
recomputation, and the recomputation is exact because its every input --- the accrued
integer, the recorded closes --- is on the record. \emph{Economic verification is re-running
the map on recorded inputs and comparing; determinism is what makes the comparison
bit-exact.}

\paragraph{(iii) The temporal residual: a safety and liveness model.}
The move algebra owns what is true \emph{within} a transaction --- conservation, coverage,
atomicity, all checked at the door. It does not own the ordering \emph{across} transactions:
which interleavings of independent firings are reachable, and whether an open obligation is
eventually discharged, are temporal questions the algebra cannot answer. The one-touch knock
while pledged is where the residual bites, and it carries an explicit model.

\medskip\noindent\textbf{Episode --- the one-touch (T2): the knock--revalue--call--discharge
model.}
W-ALPHA owns OT-1 (payout \USD{1{,}000{,}000}, marked \USD{400{,}000} pre-knock), pledged
under security-interest agreement~G to counterparty~B at a $50\%$ valuation percentage:
collateral value \USD{200{,}000} against exposure \USD{150{,}000}. OMEGA's official close
prints $79.00$, at or below the barrier $80.00$ (Chapter~\ref{ch:collateral}). Four steps
then interleave: the Monitor emits the knock; OT-1's contract triggers the unit and mints the
conditional payment obligation; the pricing layer, reading the triggered state, marks OT-1 at
$0$, so G's valuation watch opens a margin call for \USD{150{,}000}; the payment obligation
discharges \USD{1{,}000{,}000} (\USD{850{,}000} free, \USD{150{,}000} posted to~B), and the
pledge closes only after the marker returns --- an ordering the coverage invariant forces.
These steps are separate admitted transactions, and their interleaving is the temporal
residual. The model is a state machine over the reachable states, with two invariants held
always and one obligation held eventually:

\begin{lstlisting}
Coverage(s)      == \A (w,u): postedMass(w,u) =< Max(owned(w,u), 0)  -- mass; the door's invariant
NoDoubleCount(s) == (status = "triggered") => price(OT1) = 0         -- the payout is counted once
Safety   == [] ( Coverage /\ NoDoubleCount )                        -- at every committed state
Liveness == (call = "open") ~> (call = "discharged" \/ call = "closed_out")
-- Sufficiency (collateral value >= exposure) is NOT a [] invariant: it is lawfully false
-- on the interval [trigger, discharge); only Liveness constrains it. Asserting it as an
-- invariant would be a false theorem, and would forbid a lawful state.
\end{lstlisting}

\emph{Safety} says no reachable committed state double-counts value or violates coverage.
Double-counting is excluded because a triggered OT-1 prices at $0$: no committed state shows
both the option's positive mark and the \USD{1{,}000{,}000} cash it resolves into. Coverage
is excluded because it is checked at the door on every one of the interleaved transactions,
and the marker-before-extinguishment ordering of the final step is coverage refusing to
strip owned while posted mass remains. \emph{Liveness} says the margin call opened at the
revalue is eventually discharged or closed out --- the obligation-liveness requirement of the
Event Monitor made a temporal proof obligation (\S 14.1, Chapter~\ref{ch:requirements}),
under the fairness that the Events Executor eventually delivers each armed firing. The
central discipline is what the model \emph{refuses} to assert: between trigger and discharge
the book is lawfully under-collateralised, so sufficiency appears only under
$\rightsquigarrow$, never under $\square$. Coverage --- the mass guard --- and sufficiency ---
the value guard --- are the two guards of Chapter~\ref{ch:collateral}, and here the temporal
model is exactly where their difference is proved: one is a safety invariant, the other a
liveness obligation, and a model that confused them would be unsound in the direction that
either fabricates breaches or hides them.

\subsection{What an auditor checks}

A candidate implementation satisfies this chapter when the checks above run green over the
generator universe: the structural invariants of Chapter~\ref{ch:invariants} hold as
invariants and no adversarial history reaches a forbidden state; the economic oracle
reproduces every recorded transaction from its recorded event, and a seeded divergence is
repaired only by a compensating transaction; the four product-graph properties confirm that
watches, state schema, and actions never disagree; moves are shown to be the sole balance
mutator; a duplicate firing is inert and a late firing identical; firing~127 recomputes
bit-for-bit; and the knock--revalue--call--discharge model holds its two safety invariants
always and its liveness obligation eventually, with sufficiency asserted only as the latter.
Each is mechanical, each is replayable from a seed, and each fails closed. A property that
cannot be checked this way is redesigned until it can be: property-based testing is a
construction principle here, not a phase at the end.
